\documentclass[11pt]{article}
\usepackage{fontspec}
\usepackage[french]{babel}
\usepackage[leqno]{amsmath}
\usepackage{amssymb}
\usepackage{amsthm}
\usepackage{mathrsfs}
\usepackage{tikz-cd}
\usepackage{enumitem}

\newcommand{\printedpage}[1]{\marginpar{\footnotesize p.~#1}}
\newcommand{\msp}[1]{\marginpar{\footnotesize [#1]}}

% Printed arrow notations (glyphs approximated):
% Q-orthogonal couple: harpoon arrow with hooked/barbed ends
\newcommand{\Qorth}{\mathrel{\leftharpoondown\mkern-6mu\relbar\mkern-6mu\rightharpoonup}}
% the same arrow crossed by a bar
\newcommand{\Qorthx}{\mathrel{\leftharpoondown\mkern-9mu{+}\mkern-9mu\rightharpoonup}}
% harpoon tail, cross, full arrow head (as printed in a) of the first page)
\newcommand{\Qorthxfull}{\mathrel{\leftharpoondown\mkern-9mu{+}\mkern-9mu\rightarrow}}
% couple de Quillen: double arrow crossed by a bar
\newcommand{\Qcouple}{\mathrel{\leftarrow\mkern-9mu{+}\mkern-9mu\rightarrow}}
% closed couple: long double arrow
\newcommand{\clos}{\longleftrightarrow}
% long arrow with small head
\newcommand{\fl}{\longrightarrow}

\begin{document}

\printedpage{4}
% head: A. Grothendieck : Dérivateurs, Ch. XIX, Édition des Universités de Montpellier II et Paris VII
% foot: Version 11 février 2019

\msp{page 6}
\textit{alors}
\begin{enumerate}[label={}]
\item[a)] \textit{on a $\Phi_1 \Qorthxfull \Psi$ et} % reading resolved: DER_XIX_corrections.csv
\item[b)] $\Phi_1^{\sim} = \Phi^{\sim} = \Phi_1^{\natural}$.
\end{enumerate}
Ce corollaire résulte du théorème en vertu de XV prop. 4.2, dont je rappelle ici l'essentiel :

\textbf{Proposition 1.3.} \textit{Soit $\Phi \Qorthx \Psi$. Alors $\Phi^{\sim} = \Phi^{\natural} = \Psi^{*}$, et dualement on a $\Psi^{-} = \Psi^{\natural} = \Phi_{*}$.}

Notons aussi, dans cet ordre d'idées :

\textbf{Corollaire 1.4.} \textit{Soit $\Phi \Qorth \Psi$, et soit $W$ une partie quelconque de $\mathrm{Fl}(\mathcal{M})$. Pour qu'on ait $\Phi \subset W^{\natural}$, il suffit que toute flèche de $\mathcal{M}$ se factorise en $f = pi$, avec $i \in W$, $p \in \Psi$. (Et cette condition est évidemment aussi nécessaire si $\Phi \Qorthx \Psi$, et $W$ stable par facteurs directs.)}

On dit qu'une partie $\Phi_0$ de $\Phi$ est Q-\textit{génératrice à gauche} de $\Phi$ si elle satisfait la condition
\[
\Phi \ \subset \Phi_0^{\sim}\ , \tag*{(1.9)}
\]
et on dit que $\Phi$ est une partie \textit{spéciale} si $\Phi$ admet une \textit{petite} partie Q-génératrice à gauche. Si $\Phi_0 \subset \Phi$ est petit, et si $\mathcal{M}$ satisfait les conditions du théorème 1.1, alors $\Phi_0$ est une partie Q-génératrice (à gauche) si et seulement si son Q-saturé total (à gauche) $\Phi_0'''$ contient $\Phi$ (condition en apparence plus forte que $\Phi \subset \Phi_0^{\sim}$, puisque en général on a $\Phi_0''' \subset \Phi_0^{\sim}$). Cela signifie donc aussi, via 1.2, que tout \msp{page 7} élément de $\Phi$ est facteur direct d'un composé transfini de flèches déduites de
\[
\Phi_0 \cup (\text{Identités de } \mathcal{M})
\]
par cochangement de base.

Je rappelle encore le sorite $\pm$ tautologique des foncteurs adjoints :

\textbf{Proposition 1.5.} \textit{Soit}
\[
\begin{tikzcd}
\mathcal{M}' \arrow[r, shift right=1, "v^{*}"'] & \mathcal{M} \arrow[l, shift right=1, "v_{!}"']
\end{tikzcd}
\tag*{(1.5.1)}
\]
\textit{un couple de foncteurs adjoints, soient}
\[
\Phi \ \subset\ \mathrm{Fl}(\mathcal{M})\ , \qquad \Psi' \ \subset\ \mathrm{Fl}(\mathcal{M}')\ ,
\]
\textit{alors on a}
\[
(v_{!}(\Phi) \Qorth \Psi') \qquad \Longleftrightarrow \qquad (\Phi \Qorth v^{*}(\Psi'))\ . \tag*{(1.5.2)}
\]

\textbf{Corollaire 1.6.} \textit{Pour toute partie $\Phi$ de} $\mathrm{Fl}(\mathcal{M})$,
\[
(v_{!}(\Phi))_{*} \ =\ (v^{*})^{-1}(\Phi_{*})\ . \tag*{(1.6.1)}
\]

\printedpage{5}
% head: A. Grothendieck : Dérivateurs, Ch. XIX, Édition des Universités de Montpellier II et Paris VII
% foot: Version 11 février 2019

En effet, appliquant 1.5 à $\Phi$ et à $\{f'\}$, où $f' \in \mathrm{Fl}(\mathcal{M}')$, on trouve
\[
f' \in (v_{!}(\Phi))_{*} \qquad \Longleftrightarrow \qquad v^{*}(f') \in \Phi_{*}\ ,
\]
ce qui équivaut à (1.6.1). Dualement :

\textbf{Corollaire 1.7.} \textit{Pour toute partie $\Psi'$ de $\mathcal{M}'$, on a}
\[
(v^{*}(\Psi'))^{*} \ =\ (v_{!})^{-1}(\Psi'^{*}). \tag*{(1.7.1)}
\]

En effet, si $f \in \mathrm{Fl}(\mathcal{M})$, on a par (1.5.2)
\[
f \in (v^{*}(\Psi'))^{*} \qquad \Longleftrightarrow \qquad v_{!}(f) \in \Psi^{*}\ . % [sic] \Psi^{*} -> \Psi'^{*}, cf. (1.7.1)
\]
\textsuperscript{(}\footnote[3]{donner ici 1.9 page 9.}\textsuperscript{)}.

\msp{page 8}
Un couple Q-orthogonal $\Phi \Qorth \Psi$ est dit \textit{spécial} s'il satisfait les conditions suivantes.
\[
\left\{
\begin{array}{ll}
\text{a)} & \parbox[t]{0.7\textwidth}{$\mathcal{M}$ est stable par petites $\varinjlim$ , les objets de $\mathcal{M}$ sont accessibles. (On dit alors que $\mathcal{M}$ est “spéciale”.)} \\[3ex]
\text{b)} & \text{Le couple est clos : } \Phi \clos \Psi. \\[1ex]
\text{c)} & \Phi \text{ est \textit{spéciale.}}
\end{array}
\right.
\tag*{(1.10)}
\]
Cela implique donc, par 1.1, que $\Phi \Qcouple \Psi$ est un couple de Quillen. Si $\mathcal{M}$ satisfait a), les couples Q-orthogonaux spéciaux dans $\mathcal{M}$, ou encore les couples de Quillen spéciaux, correspondent donc biunivoquement aux parties $\Phi$ de $\mathrm{Fl}(\mathcal{M})$ qui sont \textit{spéciales} et closes à gauche (ou ce qui revient au même, Q-saturées à gauche), en associant à une telle $\Phi$ le couple $\Phi \Qcouple \Phi_{*}$.

Ceci posé :

\marginpar{$\Vert$}% double vertical bar printed in the left margin beside Corollaire 1.8
\textbf{Corollaire 1.8.} \textit{Supposons que $\mathcal{M}$, $\mathcal{M}'$ satisfassent la condition} a) \textit{de} (1.10), \textit{et soit $(\Phi, \Psi)$ un couple de Quillen spécial dans $\mathcal{M}$. Alors $\Bigl(\Phi' = v_{!}(\Phi)^{\sim}, \Psi' = (v^{*})^{-1}(\Psi)\Bigr)$ est un couple de Quillen spécial dans $\mathcal{M}'$.}

En effet, par le corollaire 1.6 on a
\[
\Phi'_{*} \ =\ (v_{!}(\Phi))_{*} \ =\ (v^{*})^{-1}(\Psi) \ =\ \Psi'\ ,
\]
donc $\Phi' \clos \Psi'$. Il reste à prouver que $\Phi'$ est spécial. Soit $\Phi_0$ une partie petite de $\Phi$ telle que $\Phi \subset \Phi_0^{\sim}$, posons $\Phi'_0 \ =\ v_{!}(\Phi_0) \subset \Phi'$, \msp{page 9} il suffit de voir qu'on a
\[
v_{!}(\Phi) \ \subset\ v_{!}(\Phi_0)^{\sim}\ , \tag*{($*$)}
\]
d'où résultera qu'on a $\Phi' \overset{\text{déf}}{=} (v_{!}(\Phi))^{\sim} \subset v_{!}(\Phi_0)^{\sim}$. Mais on a
\[
v_{!}(\Phi_0)^{\sim} \ =\ (v_{!}(\Phi_0)_{*})^{*}\ ,
\]
\printedpage{6}
% head: A. Grothendieck : Dérivateurs, Ch. XIX, Édition des Universités de Montpellier II et Paris VII
% foot: Version 11 février 2019
et
\[
v_{!}(\Phi_0)_{*} \ =\ (v^{*})^{-1}(\underbrace{\Phi_{0\,*}}_{\Psi}) \ =\ (v^{*})^{-1}(\Psi) \ \overset{\text{déf}}{=}\ \Psi'\ ,
\]
donc
\[
v_{!}(\Phi_0)^{\sim} \ =\ \Psi'^{*} \ =\ \Phi' \ =\ v_{!}(\Phi)^{\sim}\ ,
\]
d'où ($*$). Cet argument ne fait pas appel à l'hypothèse “spéciale” sur $\mathcal{M}$, et peut s'expliciter en le

\textbf{Corollaire 1.9.} \textsuperscript{(}\footnote[4]{devrait passer avant 1.8.}\textsuperscript{)} \textit{Sous les conditions générales de} 1.5 \textit{($v_{!} : \mathcal{M} \fl \mathcal{M}'$ étant un foncteur admettant un adjoint à droite), soient $\Phi_0$, $\Phi$ deux parties de} $\mathrm{Fl}(\mathcal{M})$ \textit{telles que $\Phi \subset \Phi_0^{\sim}$ (ou ce qui revient au même, $\Phi^{\sim} \subset \Phi_0^{\sim}$). Alors $v_{!}(\Phi) \subset v_{!}(\Phi_0)^{\sim}$, i.e. $v_{!}(\Phi)^{\sim} \subset v_{!}(\Phi_0)^{\sim}$.}

En effet, cette relation équivaut aussi à $v_{!}(\Phi_0)_{*} \ \subset\ v_{!}(\Phi)_{*}$, c'est-à-dire (par 1.6) $(v^{*})^{-1}(\Phi_{0\,*}) \ \subset\ (v^{*})^{-1}(\Phi_{*})$, et est impliqué par $\Phi_{0\,*} \ \subset\ \Phi_{*}$. Cette dernière d'autre part équivaut à l'hypothèse $\Phi^{\sim} \subset \Phi_0^{\sim}$, q.e.d.

\msp{page 10}
Si $\Psi$ est une partie de $\mathrm{Fl}(\mathcal{M})$, on en déduit pour toute $I$ dans Cat une partie
\[
\left\{
\begin{array}{l}
\Psi(I) \subset \mathrm{Fl}(\mathcal{M}(I)), \quad \text{avec}\ \ \mathcal{M}(I) = \underline{\mathrm{Hom}}(I^{\circ}, \mathcal{M})\ , \\[1ex]
f \in \Psi(I) \iff f(i) \in \Psi \ \ \forall\, i \in \mathrm{Ob}\, I\ .
\end{array}
\right.
\tag*{(1.11)}
\]
Si on introduit la catégorie discrète $I_0$ sous-jacente à $\mathrm{Ob}\, I$, et le foncteur canonique
\[
v_I = v \ :\ I_0 \fl I\ , \tag*{(1.12)}
\]
on trouve un couple de foncteurs adjoints
\[
\begin{tikzcd}[row sep=small]
\mathcal{M}(I) \arrow[r, shift right=1, "v^{*}"'] & \mathcal{M}(I_0) \arrow[l, shift right=1, "v_{!}"'] \\
( & ) \arrow[l, dotted, "v_{*}"]
\end{tikzcd}
\ \simeq\ \mathcal{M}^{I_0}\ ,
\tag*{(1.13)}
\]
et on a
\[
\Psi(I) \ =\ (v^{*})^{-1}(\Psi(I_0))\ .
\]
Supposons maintenant qu'on ait
\[
\left\{
\begin{array}{rcl}
\Psi & = & \Phi_{*} \\[1ex]
\text{d'où}\ \ \Psi(I_0) & = & \Phi(I_0)_{*}
\end{array}
\right.
\tag*{(1.14)}
\]
(vérification immédiate), on conclut alors par 1.6
\[
\Psi(I) \ =\ (v_{!}(\Phi(I_0)))_{*} \ =\ (\Phi\langle I\rangle)_{*}\ , \tag*{(1.15)}
\]
avec
\[
\Phi\langle I\rangle \ \overset{\text{déf}}{=}\ v_{!}(\Phi(I_0))^{\sim} \ =\ \Psi(I)^{*}\ . \tag*{(1.16)}
\]
Si on a un foncteur
\[
u \ :\ I \fl I'\ , \tag*{(1.17)}
\]
d'où un triple de foncteurs adjoints
\printedpage{7}
% head: A. Grothendieck : Dérivateurs, Ch. XIX, Édition des Universités de Montpellier II et Paris VII
% foot: Version 11 février 2019
\msp{page 11}
\[
\left\{
\begin{tikzcd}[column sep=large]
\mathcal{M}(I) \arrow[r, shift left=3, "u_{!}"] \arrow[r, shift right=3, "u_{*}"] & \mathcal{M}(I') \arrow[l, "u^{*}"']
\end{tikzcd}
\right.
,
\tag*{(1.18)}
\]
on aura de façon $\pm$ tautologique
\[
\left\{
\begin{array}{rcl}
u_{!}(\Phi\langle I\rangle) & \subset & \Phi\langle I'\rangle \\
u^{*}(\Psi(I')) & \subset & \Psi(I)\ .
\end{array}
\right.
\tag*{(1.19)}
\]
Le corollaire 1.8 donne :

\textbf{Proposition 1.10.} \textit{Supposons que $\Phi \Qcouple \Psi$ soit un couple de Quillen} spécial \textit{dans $\mathcal{M}$, i.e. de la forme $(\Phi, \Phi_{*})$, avec $\Phi$ partie Q-close à gauche spéciale dans $\mathcal{M}$ ($\mathcal{M}$ satisfaisant la condition} (1.10. a))\textit{. Alors pour toute $I$ dans} Cat, \textit{$(\Phi\langle I\rangle, \Psi(I))$ est un couple spécial dans $\mathcal{M}(I)$, en particulier il satisfait au théorème de factorisation.}

Considérons la construction duale, qui nous fait introduire, à partir d'un couple $\Phi \clos \Psi$ dans $\mathcal{M}$, le couple
\[
\Phi(I) \ \clos\ \Psi[I]\ , \tag*{(1.20)}
\]
où $\Phi(I)$ est défini par la formule générale (1.11) (où $\Psi$ est remplacé par $\Phi$), et où $\Psi[I]$ est défini soit comme $\Phi(I)_{*}$, soit comme
\[
\Psi[I] \ =\ v_{*}(\Psi(I_0))^{-} \ =\ \Phi(I)_{*}\ . \tag*{(1.21)}
\]

\msp{page 12}
On a alors, pour $u : I \fl I'$
\[
\left\{
\begin{array}{rcl}
u^{*}(\Phi(I')) & \subset & \Phi(I) \\
u_{*}(\Psi[I]) & \subset & \Psi[I']\ .
\end{array}
\right.
\tag*{(1.22)}
\]
Il se pose encore la question si le couple $(\Phi(I), \Psi[I])$ satisfait au théorème de factorisation, sous des hypothèses convenables sur $(\Phi, \Psi)$. Je peux seulement hasarder la

\textbf{Question-conjecture 1.11.} \textit{Supposons que $(\Phi, \Psi)$ soit un couple de Quillen} spécial, \textit{en est-il de même de $(\Phi(I), \Psi[I])$ pour toute $I$ dans} (Cat) \textit{? En d'autres termes, si $\mathcal{M}$ est spéciale, condition} (1.10. a))\textit{, et si $\Phi$ est une partie Q-close à gauche} spéciale \textit{de $\mathcal{M}$, i.e. admettant une petite partie Q-génératrice, est-il vrai que pour toute $I$ dans} Cat, \textit{$\Phi(I)$ est également spéciale ? Je conjecture que oui, en faisant au besoin des hypothèses supplémentaires sur $\mathcal{M}$ (existence d'une filtration cardinale ?). P.ex. si $\Phi_0$ est Q-génératrice dans $\Phi$, et petite, alors $\Phi_0(I)$ est-elle Q-génératrice dans $\Phi(I)$ ?}

\textbf{Proposition 1.12} \textsuperscript{(}\footnote[5]{cf. XVII 6.4 (p. 117), mais où les notations pour $\Phi\langle I\rangle$, $\Psi(I)$ sont inverses.}\textsuperscript{)}. \textit{Soit $\Phi \clos \Psi$ dans $\mathcal{M}$, soit $u : I \fl I'$ une flèche dans} Cat, \textit{d'où}
\[
\left\{
\begin{tikzcd}[column sep=large]
\mathcal{M}(I) \arrow[r, shift left=3, "u_{!}"] \arrow[r, shift right=3, "u_{*}"] & \mathcal{M}(I') \arrow[l, "u^{*}"']
\end{tikzcd}
\right.
.
\tag*{(1.12.1)}
\]
\printedpage{8}
% head: A. Grothendieck : Dérivateurs, Ch. XIX, Édition des Universités de Montpellier II et Paris VII
% foot: Version 11 février 2019
\textit{On a alors} (1.19), \textit{i.e.}
\[
\left\{
\begin{array}{rcl}
u_{!}(\Phi\langle I\rangle) & \subset & \Phi\langle I'\rangle \\
u^{*}(\Psi(I')) & \subset & \Psi(I)\ ,
\end{array}
\right.
\tag*{(1.13.1)}
\]
\textit{et de plus l'équivalence}
\[
u^{*}(\Phi(I')) \subset \Phi\langle I'\rangle \qquad \Longleftrightarrow \qquad u_{*}(\Psi(I)) \subset \Psi(I')\ . \tag*{(1.13.2)}
\]
\textit{Ces relations équivalentes sont satisfaites dans chacun des trois cas}~\textsuperscript{(}\footnote[6]{\textbf{N.B.} On a $1^{\circ} \Longrightarrow 3^{\circ}$, $2^{\circ} \Longrightarrow 3^{\circ}$.}\textsuperscript{)} \textit{:}
\begin{enumerate}[label={}]
\item[1°)] \textit{$u$ est coétale.}
\item[2°)] \textit{$I$ est discrète.}
\item[3°)] $\forall\, i' \in \mathrm{Ob}\, I'$\textit{, toute composante connexe de $I/i'$ admet un objet final.}
\end{enumerate}
\msp{page 13}

\section*{2\quad Rappels sur les triples de Quillen clos et les triples de Quillen spéciaux}

Soit $(\mathcal{M}, W)$ une \textit{catégorie de modèles,} i.e. $\mathcal{M}$ munie d'un localiseur $W \subset \mathrm{Fl}(\mathcal{M})$, i.e. une partie satisfaisant
\[
\parbox{0.75\textwidth}{$W$ contient les isomorphismes, et si $f$, $g$ sont des flèches composables, si deux parmi $f$, $g$, $gf$ sont dans $W$, la troisième aussi \textit{(saturation faible).}}
\tag*{(2.1)}
\]
Je rappelle que la donnée d'un \textit{triple de Quillen clos} $(W, C, F)$\textit{,} ou d'une structure de Quillen close, sur $\mathcal{M}$, ayant $W$ comme ensemble de quasi-isomorphismes, revient à la donnée de deux couples de Quillen
\[
C \Qcouple TF\ , \qquad TC \ \Qcouple\ F\ , \tag*{(2.2)}
\]
satisfaisant
\[
TC \ \subset\ C \cap W\ , \qquad TF \ \subset\ F \cap W\ , \tag*{(2.3)}
\]
plus une des trois conditions équivalentes :
\[
\left\{
\begin{array}{lrcll}
\text{a)} & TC & = & C \cap W & \qquad (\text{i.e. } C \cap W \subset TC) \\[1ex]
\text{b)} & TF & = & F \cap W & \qquad (\text{i.e. } F \cap W \subset TF) \\[1ex]
\text{c)} & W & = & TF \circ TC\ . &
\end{array}
\right.
\tag*{(2.4)}
\]
Rappelons aussi que ceci implique que $W$ \textit{est fortement saturée,} et a fortiori que $W = W^{\natural}$, i.e. $W$ stable par facteurs directs.

Rappelons aussi qu'en vertu de (1.4), on a :

\end{document}
